




\vspace{-10mm}
\section{Categories of Adversarial Attacks}
\label{sec:attack_types}
\vspace{-0.5em}

We group adversarial attacks into three macro classes—\textbf{Jailbreak}, \textbf{Control Generation}, and \textbf{Performance Degradation}—each revealing a distinct axis of alignment failure: ethical, semantic, and functional. 


\textbf{Jailbreak Attacks} explicitly bypass safety constraints to elicit unsafe content. These include (a) \emph{optimization-based prompts} targeting societal harm, privacy leakage, or disinformation~\cite{wu2024llms, pair23, tap23}, and (b) \emph{long-tail exploits} that trigger unsafe outputs via rare phrasing or manipulative edge cases~\cite{jiang2023promptpacker, schulhoff2023hackaprompt}.

\textbf{Control Generation Attacks} erode controllability. (a) \emph{Direct} variants involve syntax perturbations or malicious suffixes~\cite{jiang2023promptpacker}, while (b) \emph{indirect} forms hijack conditioning via goal drift~\cite{chen2024pseudo}, prompt leakage~\cite{li2024pleak}, or adversarial retrieval from external content~\cite{greshake2023indirect}.

\textbf{Performance Degradation Attacks} reduce model reliability without triggering overt refusal. These include (a) \emph{dataset poisoning} causing label flipping or semantic drift~\cite{greshake2023indirect}, and (b) \emph{prompt-based degradation} in factuality or consistency~\cite{greshake2023indirect}.





\vspace{-2mm}
\section{Too Many Attacks, Too Few Defenses}
\vspace{-0.5em}

Despite mounting evidence of alignment vulnerabilities, defenses against adversarial threats remain fractured and brittle. As attacks evolve—from prompt-level manipulations to embedding-space perturbations—they increasingly bypass safety filters not by brute force, but by exploiting structural blind spots. Most defenses remain reactive, targeting surface symptoms rather than the underlying representational geometry.


\vspace{-1mm}
\begin{table}[ht!]
\centering
\caption{\textbf{Defense Strategies Against Adversarial Attacks in LLMs.} Overview of defense paradigms, core methods, and structural limitations. Robustness remains a structurally distinct problem from alignment.}
\vspace{-1mm}
\large
\resizebox{\columnwidth}{!}{%
\begin{tabular}{@{}p{2.9cm}p{4.5cm}p{5.4cm}c@{}}
\toprule
\textbf{Defense} & \textbf{Representative} & \textbf{Limitations} & \textbf{Scalable \&} \\
\textbf{Class} & \textbf{Methods} & & \textbf{Generalizable} \\
\midrule
\textbf{Prompt-Level} & Perplexity filtering~\cite{jain2023baseline}, adversarial paraphrasing~\cite{phute2023jailbreak}, BPE-dropout & Surface-level; brittle under paraphrase or multi-hop jailbreaks & \ding{55} \\
\textbf{Training-Time} & Embedding perturbation~\cite{xhonneux2024robustness}, latent adversarial regularization~\cite{sheshadri2024latent} & High compute cost; objective- and task-sensitive & \ding{55} \\
\textbf{Certified} & Erase-and-Check~\cite{kumar2023certifying} & Narrow coverage; limited scalability and generality & \ding{55} \\
\textbf{Inference-Time} & Rewindable decoding (RAIN~\cite{li2024rain}), auxiliary vetoing~\cite{phute2023jailbreak} & Runtime overhead; dependence on auxiliary agents & \ding{55} \\
\textbf{Latent-Space} & Activation monitoring~\cite{templeton2024activations}, circuit rerouting (Cygnet~\cite{zou2024cygnet}) & Fragile under shift; depends on subspace identification & \ding{55} \\
\textbf{Geometric Alignment (Ours)} & \textbf{GRACE} (this paper) & Modular, architecture-agnostic supervision; avoids decoder modification & \ding{51} \\
\bottomrule
\end{tabular}
}
\vspace{-6mm}
\label{tab:llm_defense_taxonomy}
\end{table}


Crucially, \textit{alignment is not robustness}. Alignment governs desirable behavior under cooperative prompts; robustness demands invariance under adversarial optimization~\cite{jain2023baseline, chen2023jailbreaker}. Most defenses fail because they conflate alignment with robustness—addressing surface-level artifacts while overlooking structural vulnerabilities \textit{across the model stack} (see Table~\ref{tab:llm_defense_taxonomy}).


%%%%%%%%%%%%%%%%%%%%%%%%%%%%%%%%%%%%%%%%%%%%%%%%%%%%%%%%%%%%%%%%%%%%%%%%%%

\begin{figure*}[ht!]
    \vspace{-3mm}
    \centering
    \includegraphics[width=\textwidth]{figures/post_asr.png}
    \vspace{-6mm}
    \caption{
    \textbf{GRACE Mitigation Performance Across Open-Source LLMs.}  
    This heatmap reports \textbf{Attack Success Rate (ASR)} across 17 open-source LLMs and 12 adversarial attack types. For each attack, we show both \textbf{pre-} and \textbf{post-GRACE} ASR, with post-GRACE rows outlined in \textcolor{gold}{gold}. Each cell displays the updated ASR (rounded) and relative reduction (\%) in a two-line format.  
    \textbf{GRACE} consistently lowers ASR across diverse architectures—including instruction-tuned and chat-optimized models like \texttt{Llama-2/3}, \texttt{Vicuna}, \texttt{Mistral}, \texttt{Gemma}, and \texttt{DeepSeek}—without task-specific finetuning.  
    Attacks such as \textsc{Goal Hijacking}, \textsc{Prompt Extraction}, and \textsc{TAP} show marked mitigation, underscoring GRACE’s strength against structural and semantic adversaries. This benchmark affirms GRACE as a \textbf{robust}, \textbf{generalizable}, and \textbf{usable} safety alignment method.
    }
    \label{fig:grace_benchmark}
    \vspace{-6mm}
\end{figure*}









%%%%%%%%%%%%%%%%%%%%%%%%%%%%%%%%%%%%%%%%%%%%%%%%%%%%%%%%%%%%%%%%%%%%%%%%%%


%\textbf{Our Contribution.} We propose \textbf{GRACE}—a paradigm for LLM safety that reframes adversarial safety as a latent-space property. GRACE enforces structural constraints: \textbf{(i)} safe and unsafe completions become linearly separable; \textbf{(ii)} adversarial generations collapse into a low-entropy, isolatable submanifold. By reconfiguring representational geometry, GRACE transforms latent space into a first-class defense layer—shifting the burden from post-hoc refusal to proactive separation.



\begin{comment}
\textbf{Prompt-Level Defenses.} Surface-layer techniques such as perplexity filtering \cite{jain2023baseline}, adversarial paraphrasing \cite{phute2023jailbreak}, and BPE-dropout inject randomness to disrupt brittle suffixes, but falter against adaptive attacks.

\textbf{Training-Time Defenses.} Embedding-space perturbation \cite{xhonneux2024robustness} and latent adversarial regularization \cite{sheshadri2024latent} move the battleground deeper into the model’s computation, mitigating failure trajectories—but at high computational cost.

\textbf{Certified Defenses.} Erase-and-Check \cite{kumar2023certifying} masks and verifies substrings to yield provable robustness bounds, yet its scalability and scope remain limited.

\textbf{Inference-Time Defenses.} Dynamic safeguards like rewindable decoding (e.g., RAIN \cite{li2024rain}) and auxiliary self-vetoing models \cite{phute2023jailbreak} offer runtime flexibility, but increase latency and trust dependencies.

\textbf{Latent-Space Defenses.} Activation monitoring \cite{templeton2024activations} and circuit-based rerouting \cite{zou2024cygnet} target the representational origin of misalignment, yet depend on identifying and covering adversarial subspaces precisely.
\end{comment}




